%
% 42-cauchyliouville.tex -- Die Cauchy-Ungleichungen und der Satz von Liouville
%
% (c) 2026 Prof Dr Andreas Müller
%
\subsection{Die Cauchy-Ungleichungen und der Satz von Liouville
\label{buch:analytisch:ganz:subsection:cauchy-ungleichung}}
Die Cauchy-Integralformel für die Koeffizienten der Potenzreihenentwicklung
einer Funktion schränkt ein, wie gross die Koeffizienten werden
können.

\begin{satz}[Cauchy-Ungleichung]
\label{buch:analytisch:ganz:satz:cauchy-ungleichung}
\index{Cauchy-Ungleichung}%
Sei $f$ ein holomorphe Funktion, die auf einem offenen Gebiet definiert
ist, welches die abgeschlossene Kreisscheibe vom Radius $r$ um den
Nullpunkt enthält.
Sei ausserdem $M$ so, dass $|f(z)| \le M$ für alle $z$ mit $|z|=r$.
Dann gilt für die Koeffizienten $c_k$ der Potenzreihenentwicklung
die \emph{Cauchy-Ungleichung}
\begin{equation}
|c_k| \le \frac{M}{r^k}
\label{buch:analytisch:ganz:satz:cauchy-ungleichung}
\end{equation}
für alle $k\in\mathbb{R}$.
\end{satz}

\begin{proof}
Sei $\gamma$ der Kreis vom Radius $r$ um den Nullpunkt.
Aus der Formel
\begin{align*}
c_k
&=
\frac{1}{2\pi i} \oint_\gamma
\frac{f(z)}{z^{k+1}}\,dz
\intertext{folgt}
|c_k|
&=
\biggl|
\frac{1}{2\pi i} \oint_\gamma
\frac{f(z)}{z^{k+1}}\,dz
\biggr|
\\
&\le
\frac{1}{2\pi}
\int_0^1
\frac{|f(\gamma(t))|}{|\gamma(t)|^{k+1}}
\cdot
|\dot{\gamma}(t)|
\,dt
\\
&\le
\frac{1}{2\pi}
\cdot
\frac{M}{r^{k+1}}
\int_0^1
|\dot{\gamma}(t)|
\,dt.
\intertext{Das verbleibende Integral berechnet die Länge des
Weges $\gamma$.
Da $\gamma$ ein Kreis vom Radius $r$ ist, ist der Wert des
Integrals $2\pi r$.
Somit wird die rechte Seite zu}
&=
\frac{1}{2\pi}
\cdot
\frac{M}{r^{k+1}}
\cdot
2\pi r
=
\frac{M}{r^k}.
\end{align*}
Dies ist die behauptete Ungleichung.
\end{proof}

\begin{satz}[Liouville]
\label{buch:analytisch:ganz:satz:liouville}
\index{Satz!von Liouville}%
Sei $f(z)$ eine ganze Funktion und seien $a\in\mathbb{R}$ und
$N\in\mathbb{N}$ Zahlen mit der Eigenschaft, dass
\[
|f(z)| \le a(1+|z|^N)
\]
für alle $z\in\mathbb{C}$, dann ist $f(z)$ ein Polynom vom Grad nicht
grösser als $N$.
\end{satz}

\begin{proof}
Die Voraussetzungen des Satzes sagen, dass auf einem Kreis
vom Radius $r$ die Ungleichung
\begin{equation}
|f(z)|
\le
a(1+r^N)
\label{buch:analytisch:ganz:liouville:proof1}
\end{equation}
gilt.

Der Konvergenzradius der Potenzreihe einer ganzen Funktion ist
unendlich gross, insbesondere kann man die cauchysche Ungleichung
\eqref{buch:analytisch:ganz:satz:cauchy-ungleichung}
für beliebig
grosse Radien anwenden.
Dabei ist für die Schranke $M$ die rechte Seite der Ungleichung
\eqref{buch:analytisch:ganz:liouville:proof1}
zu verwenden.
Der Koeffizient vom Grad $n$ der Potenzreihe erfüllt damit
\begin{align*}
|c_n|
&\le 
\frac{M}{r^n}
=
a\frac{1+r^N}{r^n}
\end{align*}
für alle $r>0$.
Wir betrachten den Fall $n>N$ und lassen $r\to\infty$ gehen.
In
\begin{align*}
\lim_{r\to\infty}
a\frac{1+r^N}{r^n}
&=
a
\lim_{r\to\infty} \frac{1}{r^n}
+
a
\lim_{r\to\infty} \frac{1}{r^{n-N}}
\intertext{verschwinden beide Grenzwerte wegen $n-N>0$ auf der rechten
Seite.
Somit folgt}
|c_n|
&
=0.
\end{align*}
Somit verschwinden alle Koeffizienten der Potenzreihe für $n>N$,
die Potenzreihe ist daher nur ein Polynom vom Grad $N$.
\end{proof}

Ein wichtiger Spezialfall ist $N=0$, der oft ebenfalls der
Satz von Liouville genannt wird.

\begin{korollar}[Liouville]
Ist $f$ eine beschränkte ganze Funktion, dann ist $f$ konstant.
\end{korollar}

\begin{beispiel}
Die Funktionen $\sin z$ und $\cos z$ sind auf der reellen
Achse, aber nicht konstant.
Wären die Funktionen für alle Argumente beschränkt, müssten sie
nach dem Satz von Liouville konstant sein.
Somit müssen $\sin z$ und $\cos z$ für gewisse Argumente unbeschränkt
anwachsen.
\end{beispiel}

